%% ══════════════════════════════════════════════════════════════════════
%% NLP Corpus Analysis — Group Report
%% Based directly on acl_latex.tex (ACL 2023 official template)
%%
%% BEFORE SUBMITTING:
%%   Change \usepackage[review]{acl}  →  \usepackage{acl}
%%   Delete all % TODO lines
%%
%% FILES NEEDED IN OVERLEAF PROJECT ROOT:
%%   main.tex  (this file)
%%   custom.bib
%%   acl.sty
%%   acl_natbib.bst
%%   figures/  (folder — upload your images here)
%%
%% ── PAGE BUDGET ──────────────────────────────────────────────────────
%% Rubric: Abstract 0.25p | Intro 1.5p | Methods 3.5p | Eval 3.5p | Concl 1p
%%
%% Figure space consumed (estimate — adjust as you add/remove figures):
%%   \begin{figure}   (single column, ~5cm tall)  ≈ 0.20 pages
%%   \begin{figure*}  (full width,   ~8cm tall)   ≈ 0.38 pages
%%   \begin{table}    (small)                      ≈ 0.15 pages
%%
%% Target words per member (assuming ~2 pages of figures/tables total):
%%   Available text pages ≈ 8 pages × 650 words/page ≈ 5,200 words
%%   Per member: ~1,000 words  (Member 5 also writes Abstract: ~160w extra)
%%
%% ── FIGURE RULES (read before inserting any figure) ─────────────────
%%   • Minimum font inside figure: 9pt
%%   • Minimum resolution: 200 dpi  (save as PDF for best quality)
%%   • Single-column figure: use {figure},  width=\columnwidth
%%   • Full-width figure:    use {figure*}, width=\linewidth
%%   • Two side-by-side:     width=0.48\linewidth each, \hfill between
%%
%%   Template — single column:
%%     \begin{figure}[t]
%%       \includegraphics[width=\columnwidth]{figures/YOUR_FILE.pdf}
%%       \caption{Caption text.}
%%       \label{fig:YOUR_LABEL}
%%     \end{figure}
%%
%%   Template — full width:
%%     \begin{figure*}[t]
%%       \includegraphics[width=\linewidth]{figures/YOUR_FILE.pdf}
%%       \caption{Caption text.}
%%       \label{fig:YOUR_LABEL}
%%     \end{figure*}
%%
%%   Template — two side by side (full width):
%%     \begin{figure*}[t]
%%       \includegraphics[width=0.48\linewidth]{figures/LEFT.pdf} \hfill
%%       \includegraphics[width=0.48\linewidth]{figures/RIGHT.pdf}
%%       \caption{Caption text.}
%%       \label{fig:YOUR_LABEL}
%%     \end{figure*}
%%
%% ── TABLE RULES ─────────────────────────────────────────────────────
%%   Template:
%%     \begin{table}[t]
%%       \centering\small
%%       \begin{tabular}{lcc}
%%         \hline
%%         \textbf{Col A} & \textbf{Col B} & \textbf{Col C} \\
%%         \hline
%%         row 1          & val            & val            \\
%%         \hline
%%       \end{tabular}
%%       \caption{Caption text.}
%%       \label{tab:YOUR_LABEL}
%%     \end{table}
%%
%% ── CITATION RULES ──────────────────────────────────────────────────
%%   Add BibTeX entries to custom.bib, then:
%%     \citep{key}   →  (Author, Year)   parenthetical
%%     \citet{key}   →  Author (Year)    textual
%%   References section does NOT count toward the 10-page limit.
%% ══════════════════════════════════════════════════════════════════════

\documentclass[11pt]{article}

% Change "review" to "acl" (remove [review]) before final submission
\usepackage[review]{acl}

% Standard packages from ACL template — do not change
\usepackage{times}
\usepackage{latexsym}
\usepackage[T1]{fontenc}
\usepackage[utf8]{inputenc}
\usepackage{microtype}
\usepackage{inconsolata}
\usepackage{graphicx}

% Additional packages for this report
\usepackage{amsmath}    % for equations
\usepackage{booktabs}   % optional: \toprule \midrule \bottomrule (nicer tables)

% Set path for figures
\graphicspath{{figures/}}

% ── Helper command: delete before submitting ─────────────────────────
\usepackage{xcolor}
\newcommand{\TODO}[2]{%
  \par\noindent\colorbox{yellow!30}{\parbox{\linewidth}{%
    \textbf{[#1]}\quad\textit{\small Target: #2 words}}}%
  \par\medskip}

% ══════════════════════════════════════════════════════════════════════
\title{How Has NLP Research Changed? A Comparative Corpus Analysis
       of arXiv Abstracts, 1991--2021}

% Use \And to separate authors on the same line.
% Use \AND to force a line break between rows of authors.
\author{Member 1 \And Member 2 \And Member 3 \And Member 4 \And Member 5 \\
        Text Analytics Coursework}

\begin{document}
\maketitle

% ══════════════════════════════════════════════════════════════════════
% ABSTRACT   ~160 words | Owner: Member 5 | Write LAST
% ══════════════════════════════════════════════════════════════════════
\begin{abstract}
\TODO{Member 5 — write after all sections are finalised.
State: (1) the research question; (2) corpus and scale in one sentence;
(3) the two analytic axes and alternatives tested in two sentences;
(4) the key finding that all (or most) methods agree upon.
Maximum 160 words. Must not be a copy of the Introduction.}{160}
\end{abstract}


% ══════════════════════════════════════════════════════════════════════
% 1. INTRODUCTION   Total ~800 words | ~200 words each
%    Rubric: task, motivations, available data, requirements, examples
% ══════════════════════════════════════════════════════════════════════
\section{Introduction}
\label{sec:intro}

% Member 1 — Task and motivation (~200w) ─────────────────────────────
% Why does studying the evolution of NLP research framing matter?
% End with the explicit research question this project answers.
\TODO{Member 1 — Task and motivation.
Explain why tracking how NLP research is framed and communicated
matters scientifically. End the paragraph with a clear statement
of the research question.}{200}

% Member 2 — Available data (~200w) ──────────────────────────────────
% Describe the corpus: source, size, time range, period binning,
% notable characteristics. Add a figure or table here if useful.
\TODO{Member 2 — Available data.
Describe the arXiv cs.CL corpus: where it comes from, total N,
time range (1991--2021), how it is split into 5-year periods,
and any notable characteristics such as class imbalance.
Insert a figure or table if it helps (use the templates in the header).}{200}

% Member 3 — Requirements for a good solution (~200w) ─────────────────
% What must a method do well to answer the research question?
% Use these requirements to motivate the two analytic axes.
\TODO{Member 3 — Requirements for a good solution.
List 3--4 properties a method must have to answer the research
question reliably. Use them to motivate the two analytic axes
introduced in Section~\ref{sec:methods}.}{200}

% Member 4 — Dataset examples (~200w) ────────────────────────────────
% Paraphrase one early and one recent abstract to illustrate the shift.
% Add a data-exploration figure if you think it strengthens the section.
\TODO{Member 4 — Dataset examples.
Paraphrase one early (pre-2005) and one recent (post-2015) abstract.
Identify 2--3 concrete differences in vocabulary and framing.
A data-exploration figure is optional but encouraged.}{200}


% ══════════════════════════════════════════════════════════════════════
% 2. METHODS   Total ~1,800 words | ~360 words each
%    Rubric 3a: baseline system, preprocessing, representations
%    Rubric 3b: axes of experimentation
%    Rubric 3c: proposed improvements + hypotheses
% ══════════════════════════════════════════════════════════════════════
\section{Methods}
\label{sec:methods}

% Member 1 — Preprocessing + TF-IDF baseline (Rubric 3a) (~360w) ─────
% Describe the preprocessing pipeline and the TF-IDF baseline.
% Include the TF-IDF equation. State vocabulary sizes. Justify choices.
% Add a pipeline diagram or figure if it makes the design clearer.
\TODO{Member 1 — Preprocessing and TF-IDF baseline (Rubric 3a).
Cover: overall pipeline in brief, preprocessing steps with justifications,
TF-IDF unigram baseline including the equation below, vocabulary size and
hyperparameters, TF-IDF bigram as the first Axis-1 alternative.}{360}

\begin{equation}
  \mathrm{TF\text{-}IDF}(t,d) =
    \mathrm{tf}(t,d)\cdot\log\!\frac{N}{\bigl|\{d':t\in d'\}\bigr|}
  \label{eq:tfidf}
\end{equation}

% Member 2 — SBERT + LDA (Rubric 3a + 3c) (~360w) ─────────────────────
% Describe SBERT and LDA as proposed Axis-1 improvements.
% Justify the model/configuration choices. State a hypothesis for each.
\TODO{Member 2 — SBERT and LDA representations (Rubric 3a + 3c).
Explain SBERT \citep{reimers-2019-sentence-bert}: model used, output
dimensionality, why it captures semantic drift beyond vocabulary overlap.
Explain LDA \citep{blei2003latent}: number of topics K, how K was chosen,
what the topic-proportion vector represents.
State a hypothesis for each: how do you expect it to perform differently
from TF-IDF?}{360}

% Member 3 — Keyword shift axis (Rubric 3b + 3c) (~360w) ──────────────
% Describe the keyword shift analysis as Axis-2A.
% State a clear hypothesis about what it reveals.
\TODO{Member 3 — Keyword shift method and hypothesis (Rubric 3b + 3c).
Explain how top-N terms per period are extracted, how term frequency is
normalised, and what ``emerging'' vs ``disappearing'' vocabulary means.
State hypothesis: what signal do you expect keyword shifts to reveal
that other Axis-1 methods cannot?
Add a figure if it clarifies the method design.}{360}

% Member 4 — Cosine distance axis (Rubric 3b + 3c) (~360w) ────────────
% Describe cosine distance between period centroids as Axis-2B.
% State a clear hypothesis about what it reveals.
\TODO{Member 4 — Cosine distance method and hypothesis (Rubric 3b + 3c).
Explain how period centroids are computed, why cosine distance is
appropriate for high-dimensional text vectors, and how this method
differs from keyword shift.
State hypothesis: what signal does cosine distance detect that keyword
shift misses?
Add a figure if it clarifies the method design.}{360}

% Member 5 — Dimensionality reduction (Rubric 3a + 3c) (~360w) ────────
% Describe PCA and UMAP as the visualisation component.
% Justify the two-stage SVD→PCA pipeline for sparse TF-IDF matrices.
% State a hypothesis comparing PCA and UMAP.
\TODO{Member 5 — Dimensionality reduction design (Rubric 3a + 3c).
Explain the two-stage SVD(k=50)→PCA(2D) pipeline for sparse TF-IDF:
why direct PCA on sparse data is infeasible; SVD as LSA
\citep{deerwester1990indexing}; justify k.
Explain direct PCA for dense matrices (SBERT, LDA).
Explain UMAP \citep{mcinnes2018umap-software}: non-linear alternative,
pipeline for each matrix type, sensitivity to \texttt{n\_neighbors}.
State hypothesis: PCA vs UMAP---what does each reveal?
Include the explained-variance figure here if it supports k=50.}{360}


% ══════════════════════════════════════════════════════════════════════
% 3. EXPERIMENTAL EVALUATION   Total ~1,500 words | ~300 words each
%    Rubric 4a: performance metrics and limitations
%    Rubric 4b: experimental setup
%    Rubric 4c: results with interpretation
%    Rubric 4d: error analysis
%    Rubric 4e: discussion, strengths/limitations, recommendations
% ══════════════════════════════════════════════════════════════════════
\section{Experimental Evaluation}
\label{sec:eval}

% Member 1 — TF-IDF results (Rubric 4a--4d) (~300w) ──────────────────
% Report the TF-IDF trajectory findings. Add figures or tables
% you consider best for presenting these results.
% Analyse surprising or misleading results from TF-IDF.
\TODO{Member 1 — Axis-1 TF-IDF results (Rubric 4a--4d).
State the metric(s) you use for TF-IDF trajectories and their limitations.
Describe how TF-IDF hyperparameters were set.
Present your main TF-IDF findings with figures or tables as appropriate.
Identify patterns where TF-IDF produces unexpected or misleading results.}{300}

% Member 2 — SBERT + LDA results (Rubric 4a--4d) (~300w) ─────────────
% Report the SBERT and LDA trajectory findings. Add figures or tables
% as appropriate. Analyse where they agree with or diverge from TF-IDF.
\TODO{Member 2 — Axis-1 SBERT and LDA results (Rubric 4a--4d).
State the metric(s) used and their limitations for dense representations.
Describe how K (LDA) and the SBERT model were selected.
Present your main SBERT and LDA findings with figures or tables.
Where do these agree with or diverge from TF-IDF, and why?}{300}

% Member 3 — Keyword shift results (Rubric 4a--4d) (~300w) ────────────
% Report the keyword shift findings. Add a figure if helpful.
% Analyse errors: where does keyword shift mislead or fail?
\TODO{Member 3 — Axis-2A keyword shift results (Rubric 4a--4d).
State the metric(s) and their limitations.
Describe the experimental setup: N chosen, period pairs, normalisation.
Present main keyword findings with a figure if helpful.
Error analysis: generic terms? class-imbalance bias? writing-convention
artefacts in bigrams?}{300}

% Member 4 — Cosine distance results (Rubric 4a--4d) (~300w) ──────────
% Report the cosine distance findings. Add a figure if helpful.
% Analyse errors. Compare with keyword shift.
\TODO{Member 4 — Axis-2B cosine distance results (Rubric 4a--4d).
State the metric(s) and their limitations.
Describe the setup: centroid computation, distance formula, matrices used.
Present main cosine findings with a figure if helpful.
Error analysis: centroid averaging masking within-period heterogeneity?
Minimum-distance period pair? How do results compare with Member~3?}{300}

% Member 5 — Discussion and limitations (Rubric 4e) (~300w) ───────────
% Synthesise all findings. Compare Axis 1 methods. Compare Axis 2.
% Which method best balances consistency and interpretability?
% List 4--6 specific limitations. Add an evaluation table if useful.
\TODO{Member 5 — Discussion and limitations (Rubric 4e).
Cross-method comparison: do Axis-1 methods agree on the breakpoint period?
Do Axis-2 methods agree on the same breakpoint? Where they disagree, explain.
Strengths and limitations of each method (one sentence each, 4--6 methods).
An evaluation table (Consistency / Interpretability / Alignment with
known NLP history) is optional but encouraged.
End with 1--2 concrete recommendations based on your analysis.}{300}


% ══════════════════════════════════════════════════════════════════════
% 4. CONCLUSIONS   Total ~600 words | ~120 words each
%    Rubric: key findings, what worked, what did not, open challenges
% ══════════════════════════════════════════════════════════════════════
\section{Conclusions}
\label{sec:conclusions}

% Each member writes one paragraph (~120 words):
%   • Their key finding from their method/axis
%   • One open challenge they identified

% Member 5 — opening: answer the research question (~120w) ────────────
\TODO{Member 5 — Opening paragraph.
Directly answer the research question: how has NLP research changed
from 1991 to 2021? State the dominant transition period and whether
the methods agree on it.}{120}

% Member 1 — TF-IDF key finding + open challenge (~120w) ──────────────
\TODO{Member 1 — Key TF-IDF finding and one open challenge.}{120}

% Member 2 — SBERT/LDA key finding + open challenge (~120w) ──────────
\TODO{Member 2 — Key SBERT/LDA finding and one open challenge.}{120}

% Member 3 — Keyword key finding + open challenge (~120w) ─────────────
\TODO{Member 3 — Key keyword-shift finding and one open challenge.}{120}

% Member 4 — Cosine key finding + open challenge (~120w) ──────────────
\TODO{Member 4 — Key cosine-distance finding and one open challenge.}{120}

% Member 5 — closing: what worked, what did not, forward look (~120w) ─
\TODO{Member 5 — Closing paragraph.
What worked well overall. What did not work. One sentence on the most
valuable future extension.}{120}


% ══════════════════════════════════════════════════════════════════════
% REFERENCES  (not counted toward the 10-page limit)
% Add BibTeX entries to custom.bib
% ══════════════════════════════════════════════════════════════════════
\bibliography{custom}

\end{document}
